\section{Index des rapports, dérivations et explorations}\label{ann:index}
Liste générée depuis le frontmatter (\texttt{headline}) de chaque fichier, au 15 novembre 2026. Les titres sont repris tels quels : ils peuvent contenir des valeurs périmées, signalées en annexe~\ref{ann:incoherences}.
\begin{small}\begin{longtable}{@{}p{0.42\textwidth}p{0.54\textwidth}@{}}
\toprule Fichier & Titre (\texttt{headline}) \\ \midrule\endhead
\texttt{v2/docs/lab/reports/2026-11-14-accueil-personas.md} & 3 personas sur 3 refusent l'accueil actuel (France par régions, leviers ouverts, aucun chiffre avant clic) ; 3 sur 3 veulent un chiffre-phrase dès l'ouverture et les leviers repliés ; désaccord sur le point d'entrée (… \\
\texttt{v2/docs/lab/reports/2026-11-14-audit-sources.md} & 469 citations auditées : 382 confirmées, 55 corrigées, 32 non vérifiables ; aucune erreur ne renverse une conclusion, mais 3 corrections changent un chiffre affiché (quotas DD76 ±9/12/13 \%, « +1 \% par an » de retraite… \\
\texttt{v2/docs/lab/reports/2026-11-14-axe1-hyperparametres.md} & 27 fiches de paramètres : 19 sourcés, 1 mesuré, 1 supposé, 6 non mesurables ; seul effet de régulation chiffré = IDE (−25,57 IDEL / 100 000 hab., IRDES) \\
\texttt{v2/docs/lab/reports/2026-11-14-axe2-apc-sorties.md} & APC gaussien de la rétention nette : MAPE dép. h=5 8,74 \% contre 8,68 \% (A0), égalité (Δ +0,07 [−0,03 ; +0,16]) ; bat la rétention 60+ en tendance (9,00 \%) ; CRPS relatif 6,82 \% mais intervalles trop étroits (couvertu… \\
\texttt{v2/docs/lab/reports/2026-11-14-axe2-cohortes-bayesiennes.md} & Cohortes bayésiennes (APC des sorties + entrées en niveau local × parts hiérarchiques) : MAPE dép. h=5 7,73 \% contre 8,68 \% (Δ −0,95 pt, IC 95 \% [−1,51 ; −0,47]) ; bande à 80 \% couvrant 77,6 \% sur toutes les origines … \\
\texttt{v2/docs/lab/reports/2026-11-14-axe2-communes.md} & Pas encore. K (cohortes hiérarchiques commune ⊂ EPCI ⊂ dép ⊂ région) : bande trop large (couverture à 80 \% de 90–100 \% chez les IDE) et point moins bon que la part tenue sous 10 personnes ; repli K2 (part tenue × tota… \\
\texttt{v2/docs/lab/reports/2026-11-14-axe2-methodes.md} & Shrinkage hiérarchique dép → région → national : MAPE dép. h=5 8,57 \% contre 8,68 \% pour le stock-flux à taux locaux (Δ −0,11 pt, IC 95 \% [−0,21 ; −0,02], périmètre privé ; −0,10 sur l'ancien) ; âges fins à égalité ; … \\
\texttt{v2/docs/lab/reports/2026-11-14-axe3-n1-prevision.md} & Protocole pré-enregistré, pas encore lancé : critère = MAPE dép. à h = 5 ans, borne haute de l'IC 95 \% (bootstrap par département) de l'écart à la meilleure référence naïve < 0 \\
\texttt{v2/docs/lab/reports/2026-11-14-axe3-n10-grain-couverture.md} & Protocole pré-enregistré, pas encore lancé : 3 grains × 3 professions × origines 2009–2020 ; bande P10–P90 empirique par grain × classe de taille, tirée des seules erreurs passées ; « fiable » = couverture dans [75 ; … \\
\texttt{v2/docs/lab/reports/2026-11-14-axe3-n11-confirmation-k4.md} & Protocole pré-enregistré : K4 passe d'abord N10 tel quel (2014–2020, h = 5), puis une confirmation sur l'origine réservée 2021, h = 4 (cible 2025), jugée avec les mêmes critères et comparée aux références ; confirmati… \\
\texttt{v2/docs/lab/reports/2026-11-14-axe3-n12-confirmation-k5.md} & Protocole pré-enregistré, gelé avant tout jugement de K5 : même règle que N11 ; origine réservée 2022, h = 3 (cible 2025), proposée par l'axe 2, acceptée comme confirmation partielle ; références jugées sur la même or… \\
\texttt{v2/docs/lab/reports/2026-11-14-axe3-n2-long-terme.md} & Protocole pré-enregistré, pas encore lancé : reproduire le tendanciel DREES n°76 en ratios N(t)/N(2020) à ±2 \% en 2030 et ±3 \% en 2050 (CD, SF, PH), et passer trois cas limites (h → 0 exact, N* = Q·C à 1 \%, verrou du … \\
\texttt{v2/docs/lab/reports/2026-11-14-axe3-n3-leviers-did.md} & IDE : effet −30 \% [−34 ; −26] de la régulation en zone sur-dotée sur 2012–2019, mais le test pré-enregistré de tendances parallèles échoue (p = 0,043 < 0,10), donc levier NON validé selon le protocole ; kinés : la pen… \\
\texttt{v2/docs/lab/reports/2026-11-14-axe3-n4-echelle-fine.md} & Protocole pré-enregistré, pas encore lancé : N\_min = plus petite classe de taille où le P80 de l'erreur absolue à 5 ans ≤ 25 \% (3 professions, ≥ 3 origines sur 4) ; prédiction dérivée du seul bruit de Poisson : N\_min … \\
\texttt{v2/docs/lab/reports/2026-11-14-axe3-n5-mondes-simules.md} & Protocole pré-enregistré, pas encore lancé : 7 mondes simulés à vérité connue (repris de l'observatoire, A1) ; dans un monde qui respecte ses hypothèses, un estimateur doit tenir |erreur nationale| ≤ 1 \% à 15 ans ; da… \\
\texttt{v2/docs/lab/reports/2026-11-14-axe3-n6-probabiliste.md} & Protocole pré-enregistré, pas encore lancé : CRPS sur log N à h = 5 (même critère bootstrap que N1), calibration jugée sur la couverture des intervalles à 50 \% et 80 \% (bandes 44–56 \% et 75–85 \%), et calibration par s… \\
\texttt{v2/docs/lab/reports/2026-11-14-axe3-n7-addendum.md} & Protocole additionnel gelé avant calcul : (A) DiD dans les bassins de vie 2022 ; (B) verdict N1 par origine 2016, 2018 et 2021 ; (C) bandes « pente d'erreur » à armes égales pour tous les estimateurs ; (D) cible ER 13… \\
\texttt{v2/docs/lab/reports/2026-11-14-axe3-n8-bande-racine.md} & Protocole pré-enregistré, pas encore lancé : log N \textasciitilde{} Normale(log prédit, s²·h), s² = moyenne de e²/h′ sur les paires passées ; mêmes origines, critère CRPS et seuils de couverture que N7-C et N6 \\
\texttt{v2/docs/lab/reports/2026-11-14-axe3-n9-boucles.md} & Protocole pré-enregistré, pas encore lancé : B2 par bassin de vie × année (IDE, MK, CD, 2014–2025), B1 par département × année (5 professions, cohortes 45–54 → 50–59 à 5 ans) ; verdict par profession : signe confirmé,… \\
\texttt{v2/docs/lab/reports/2026-11-14-axe3-plan-validation.md} & 4 protocoles pré-enregistrés (prévision, 2050, leviers, échelle fine) gelés par empreinte SHA-256 avant tout résultat ; le niveau 3 est lancé pour les IDE \\
\texttt{v2/docs/lab/reports/2026-11-14-axe4-flux-communaux.md} & Proxy d'installations nettes par commune 2012–2025 (IDE, MK, CD ; Σ communes = département à ≤ 0,34 \% pour MK et CD) mais il ne voit pas toute l'année blanche kinés 2018 (0,90, seuil 0,8 : échec) ; 20 083 officines en… \\
\texttt{v2/docs/lab/reports/2026-11-14-axe4-grain-communal.md} & Observable à la commune pour les 5 professions (Annuaire Santé, nov. 2026 ; historique 2004–2025 pour IDE, MK, CD). Projetable : bassin de vie / EPCI pour les infirmiers (médiane 35–53 privés par unité), EPCI en hiéra… \\
\texttt{v2/docs/lab/reports/2026-11-14-axe4-projection-communale.md} & Pharmaciens par commune 2012–2026 reconstitués (Σ = DREES dép. ; erreur médiane 16 \% à la commune, 10,7 \% à l'EPCI contre l'Annuaire : test échoué de peu) ; sages-femmes : historique communal non disponible (5 routes … \\
\texttt{v2/docs/lab/reports/2026-11-14-axe4-registre-evenements.md} & 30 événements datés 2009–2028 dans v2/data/events.csv (24 vérifiés, 6 partiels, après l'audit des sources) ; l'année blanche MK 2018 se voit dans le RPPS (moins de 25 ans : 5 634 → 3 687, −34,6 \%) ; le saut d'inscrits… \\
\texttt{v2/docs/lab/reports/2026-11-14-axe4-secteur-public-prive.md} & Sages-femmes « libéral » 2026 : 12 059 – 14 660 (centrale 12 912), contre 12 475 – 26 428 ; largeur −81 \%. Pharmaciens 60 005 – 67 219, dentistes 40 784 – 47 894, kinés 98 828 – 103 986 ; infirmiers 2023 : 277 788, sa… \\
\texttt{v2/docs/lab/reports/2026-11-14-axe4-zonage-commune.md} & 34 919 géocodes × 4 professions (communes COG 2026 + 45 arrondissements PLM) ; « remplacement seul » : 97 géocodes et 4,8 \% de la population pour les dentistes (cité : 97 communes, ≈ 5 \%), 27–28 \% pour kinés et infirm… \\
\texttt{v2/docs/lab/reports/2026-11-14-axe5-sensibilite.md} & Partie 2 (jumeau b64ce56) : en 2050, la décision pèse plus que l'ignorance pour les IDE (0,89 contre 0,07), les MK (0,81 contre 0,12) et les SF (0,59 contre 0,38), à égalité pour les CD, moins pour les PH (0,38 contre… \\
\texttt{v2/docs/lab/reports/2026-11-14-curseur-annee.md} & Curseur 2026 → 2050 sans recalcul au déplacement ; /api/carte ≈ 100 ms pour 25 années, 78 ko ; 2050 inchangé \\
\texttt{v2/docs/lab/reports/2026-11-14-depot-evolution-soin.md} & Même famille de modèle que notre v1 (cohortes par tranches de 5 ans, taux de maintien). Plus avancé sur la démo et sur la validation (mondes simulés, backtest à 3–10 ans). Leur levier retraite avait un bug qui touche … \\
\texttt{v2/docs/lab/reports/2026-11-14-jumeau-axe0.md} & Version 6. 17 critères sur 17. À la commune, 5 professions en 1,35 à 1,57 s au lieu de 3,99 s (698 212 cellules, même état). Vivier réel DREES 491 : diplômés IDE reconstitués à 0,5 \% près, mais la persistance prévoit … \\
\texttt{v2/docs/lab/reports/2026-11-14-produit.md} & Pas achetable en l'état (8 conditions de signature, 0 remplie complètement) ; voie réaliste = commun numérique libre porté par une institution (DREES, ARS pilote, beta.gouv), puis prestation < 60 k€ HT par ARS (seuil … \\
\texttt{v2/docs/lab/reports/2026-11-14-prototype-carte.md} & Oui. Carte MapLibre plein écran (localhost:8767) : un mouvement de levier recolore les 101 départements et met à jour la fiche en 53 ms (p50), 66 ms (p95), 87 ms (max) de réponse perçue, n = 35, sous la cible de 150 ms \\
\texttt{v2/docs/lab/reports/2026-11-14-prototype-decision.md} & Oui. Le prototype tourne sur localhost:8766 : un scénario est recalculé en 3 à 32 ms par le jumeau et en 50 à 120 ms aller-retour. L'écart au tendanciel est décomposé exactement par flux. Exemple : kinés, quotas +20 \%… \\
\texttt{v2/docs/lab/reports/2026-11-14-reconciliation-dentistes.md} & Socle 2026 dentistes privé = 47 347 (≤ 48 690 tous modes, RPPS 2026) ; 76 397 est la projection 2050, pas le socle \\
\texttt{v2/docs/lab/reports/2026-11-14-reprise-observatoire.md} & 31 éléments passés en revue : 19 repris ou adaptés, 2 à vérifier, 7 non repris, 3 réservés à la couche personas. Les 3 plus utiles : la validation par mondes simulés, le levier retraite corrigé (+n ans exactement en 6… \\
\texttt{v2/docs/lab/reports/2026-11-14-revue-adversariale-axe0-v2.md} & Non. Les 6 constats soumis à réfutation sont confirmés, soit 4 défauts bloquants distincts : levier retraite (tranche 65+ ouverte, +41 \% à Δ = 2), régulation qui fait encore varier le total (+254 dentistes), scénario … \\
\texttt{v2/docs/lab/reports/2026-11-14-revue-adversariale-axe0-v3.md} & Oui pour les défauts bloquants : il n'en reste plus dans le jumeau. Le seul défaut bloquant était dans le prototype (plantage à la réactivation de la régulation) et il est corrigé. Restent 5 défauts importants confirm… \\
\texttt{v2/docs/lab/reports/2026-11-14-revue-adversariale-axe0.md} & Non, pas encore. 6 défauts confirmés (4 bloquants, 2 importants) sur 6 soumis à réfutation, plus 3 non vérifiés. L'écart 0,0 de la reproduction tient ; les critères plug-in, levier, agrégation et conservation sont en … \\
\texttt{v2/docs/lab/reports/2026-11-14-uat-agentique-carte.md} & 5 tâches réussies sur 21 (14 partielles, 2 échecs) ; 37 constats confirmés sur 41 soumis ; après dédoublonnage, 3 P0 (titre de fiche d'une autre profession, sélecteur de profession hors écran sur mobile, export sans s… \\
\texttt{v2/docs/lab/reports/2026-11-14-uat-carte-v2.md} & 4 tâches réussies sur 19 (12 partielles, 3 échecs) ; 30 constats confirmés sur 39 soumis ; 1 P0 (sélecteur de profession hors écran sur mobile) et 8 P1 confirmés (6 après dédoublonnage) \\
\texttt{v2/docs/lab/reports/2026-11-14-uat-interface.md} & V1 : 11 tâches réussies sur 22 (6 échecs, dont l'effectif 2035 et l'export) ; 38 constats confirmés sur 40. V2 : verdict « à porter en démo », mais recette faussée par un onglet partagé (5 tâches sur 16) ; P1 = sélect… \\
\texttt{v2/docs/lab/reports/2026-11-15-audit-rgpd-aipd-dm.md} & Pas un dispositif médical (aucune finalité médicale individuelle) ; AIPD probablement non obligatoire (2 critères G29 sur 9 au plus, données déjà publiques) ; un point RGPD réel : 52,5 \% des 1 251 610 cellules commune… \\
\texttt{v2/docs/lab/reports/2026-11-15-axe1-relecture-sources.md} & 29 paramètres relus : 20 confirmés, 8 corrigés dans leur fiche, 1 introuvable ; aucune valeur du code ne change, jumeau.check 18 sur 18 \\
\texttt{v2/docs/lab/reports/2026-11-15-budget-erreur-c1-a2.md} & MAPE C1 7,76 \% [7,12 ; 8,40] : plancher de bruit poissonien 3,11 [2,89 ; 3,36] ; l'écart local des entrées (−2,60 [2,20 ; 3,05] si connu) et du milieu de carrière (−2,36 [2,00 ; 2,71]) dominent, le national pèse 0,49 … \\
\texttt{v2/docs/lab/reports/2026-11-15-c1-gbm-global.md} & Rejeté contre A2 et contre C1 (7,73 \%, modèle servi) : le GBM global fait 10,15 \% de MAPE dép. à h = 5 contre 8,60 \% pour A2 (Δ +1,55 pt, IC 95 \% [+0,62 ; +2,75]) ; le MLP 20,52 \% (Δ +11,92, [+10,54 ; +13,86]) ; aucun… \\
\texttt{v2/docs/lab/reports/2026-11-15-c1-intervalles.md} & Non. À l'origine 2021 (test N6), aucun des 6 candidats ne bat la persistance (meilleur : C1-conf, +0,0043 [−0,0006 ; +0,0094] en CRPS sur log N), et tous sont moins bons que C1 brut (+0,0012 à +0,0098). Le conforme ab… \\
\texttt{v2/docs/lab/reports/2026-11-15-jury-grille.md} & 64/100 (fourchette 54–74) : pertinence 6,5, IA 7,0, prototype 6,5, faisabilité 5,5, pitch 6,5 ; la note de conformité (RGPD, AIPD, 2017/745) rapporte le plus de points, \textasciitilde{}3 pts pour 45 min \\
\texttt{v2/docs/lab/derivations/2026-11-14-agent-interprete.md} & Un agent à outils explique les sorties du jumeau ; il ne calcule rien, et chaque nombre de sa réponse est vérifié contre les sorties d'outils \\
\texttt{v2/docs/lab/derivations/2026-11-14-assistant-objectifs.md} & D'un objectif chiffré à un plan de leviers daté : le jumeau cherche, vérifie et date chaque mélange ; le modèle de langage traduit et explique \\
\texttt{v2/docs/lab/derivations/2026-11-14-axe2-apc-sorties.md} & — \\
\texttt{v2/docs/lab/derivations/2026-11-14-axe2-cohortes-bayesiennes.md} & — \\
\texttt{v2/docs/lab/derivations/2026-11-14-axe2-communes.md} & — \\
\texttt{v2/docs/lab/derivations/2026-11-14-axe2-familles-sd-abm.md} & — \\
\texttt{v2/docs/lab/derivations/2026-11-14-axe2-methodes.md} & — \\
\texttt{v2/docs/lab/derivations/2026-11-14-module-besoin.md} & Besoin = habitants × profil d'âge du recours ; un transfert de tâches déplace une part du besoin, il ne crée aucun professionnel \\
\texttt{v2/docs/lab/derivations/2026-11-15-c1-gbm-global.md} & — \\
\texttt{v2/docs/lab/explorations/2026-11-14-etat-de-lart-methodes.md} & Le gain viendra d'abord de meilleurs modèles statistiques sur la même structure transparente (cohortes par composantes, bayésien hiérarchique), puis de données de flux ciblées, et presque pas de nouveaux hyperparamètr… \\
\bottomrule\end{longtable}\end{small}
